\subsection*{1.2. Ordinary quadratic singularities}

Let $y$ be a closed point of a scheme $Y$ of finite type over a field
$k$ of characteristic $p$, and let $n$ be the dimension of $Y$ at
$y$.

\subsubsection*{Definition 1.2.1 (= VI 6.6).}
\begin{enumerate}[label=(\roman*),leftmargin=2.6em]
\item \underline{If $k$ is algebraically closed, $y$ is called an
ordinary quadratic point of $Y$ if the completion
$\widehat{\mathcal O}_{Y,y}$ of the local ring of $Y$ at $y$ is
isomorphic to the quotient of $k[[x_1,\ldots,x_{n+1}]]$ by the ideal
generated by a single formal power series $f$ which can be written}
\begin{equation}\tag{1.2.1.1}
f(\underline X)=Q(\underline X)+\text{terms of order }>2,
\end{equation}
\underline{where $Q(\underline X)$ is an ordinary quadratic form.}
\footnote{The printed display is numbered (4.2.1) and has an extra
equals sign after $Q(\underline X)$; the English restores the local
numbering and the asserted expansion $f=Q+$ higher-order terms.}

\item \underline{If $\bar k$ is an algebraic closure of $k$, $y$ is
called an ordinary quadratic point of $Y$ if the points of
$\bar Y=Y\otimes_k\bar k$ lying over $y$ are ordinary quadratic
points of $\bar Y$.}
\end{enumerate}

\subsubsection*{1.2.2.}
Replacing ``ordinary quadratic form'' by ``nondegenerate quadratic
form'' in this definition gives the notion of a
\underline{nondegenerate quadratic point} (VI 6.0).  A point $y$ is
a nondegenerate quadratic point of $Y$ if and only if it is an
ordinary quadratic point and either $p\ne2$ or $n$ is odd.

\subsubsection*{Example 1.2.3.}
Let $Q(\underline X)$ be a nondegenerate quadratic form in $n+1$
variables over $k$.  Let $Y_0$ be the quadratic cone in
$\mathbb A_k^{n+1}$ with equation $Q=0$.  Then the origin, denoted by
$y_0$, is a nondegenerate quadratic point of $Y_0$.

\subsubsection*{Example 1.2.4.}
Suppose that $p=2$ and $n=2m$.  Let $a\in k$, and let $Q$ be a
polynomial
\[
Q(\underline X)=(X_0^2-a)
 +\sum_{0<i\leq j\leq2m}a_{ij}x_ix_j.
\]
Suppose that the quadratic form in $2m$ variables
$\sum_{0<i\leq j}a_{ij}x_ix_j$ is nondegenerate, and denote by $y_0$
the closed point of $\mathbb A_k^{n+1}$ with coordinates
$(\sqrt a,0,\ldots,0)$.  Then $y_0$ is an ordinary quadratic point
of the closed subscheme $Y_0\subset\mathbb A_k^{n+1}$ with equation
$Q=0$.

\subsubsection*{Example 1.2.5.}
Let $k'$ be a finite separable extension of $k$, and let $Y_0$ be a
$k'$-scheme of the kind considered in 1.2.3 (respectively, 1.2.4).
If $Y_0$ is regarded as a $k$-scheme, then $y_0$ is a nondegenerate
(respectively, ordinary) quadratic point of $Y_0$.  The extension
$k'/k$ identifies with the largest separable subextension of
$k(y_0)/k$.

\subsubsection*{Theorem 1.2.6.}
\underline{Suppose that $y$ is an ordinary quadratic point of the
$k$-scheme $Y$, and let $k'$ be the largest separable subextension
of $k(y)/k$.  There then exist a $k'$-scheme
$Y_0\subset\mathbb A_{k'}^{n+1}$ of the kind in 1.2.3 or 1.2.4 and a
$k$-isomorphism $\varphi$ between the henselization $Y_{(y)}$ of $Y$
at $y$ and the henselization $Y_{0(y_0)}$ of $Y_0$ at $y_0$.}

First observe that there is an etale neighborhood $Y'$ of $y$ in
$Y$
\begin{center}
\begin{tikzcd}[row sep=large,column sep=large]
{} & Y'\arrow[d]\\
y\arrow[ur,hook]\arrow[r,hook]&Y
\end{tikzcd}
\end{center}
which is a $k'$-scheme.  Indeed, $Y_{(y)}$ is canonically a
$k'$-scheme.  Replacing $(k,Y,y)$ by $(k',Y',y)$ reduces us to the
case $k=k'$, that is, to the case in which $k(y)/k$ is radicial.
